\documentclass[11pt,a4paper]{article}
\usepackage[italian]{babel}
\usepackage[utf8]{inputenc}
\usepackage[T1]{fontenc}
\usepackage{amsmath,amssymb,amsthm}
\usepackage{geometry}
\geometry{margin=2.3cm}
\usepackage{enumitem}
\usepackage{tikz}
\usetikzlibrary{angles,quotes,arrows.meta}
\usepackage{tabularx}
\usepackage{xcolor}
\usepackage{hyperref}
\hypersetup{colorlinks=true, linkcolor=blue, urlcolor=blue}

% ---------- Ambiente "Esempio" ----------
\newcounter{esempio}[section]
\newenvironment{esempio}[1][]{%
  \refstepcounter{esempio}%
  \par\vspace{4pt}\noindent
  \textbf{\color{blue!60!black}Esempio \theesempio.} #1\ \par\nobreak
  \begingroup\color{black}
}{%
  \endgroup\par\vspace{4pt}
}

% ---------- Ambiente "Osservazione" ----------
\newenvironment{osservazione}{%
  \par\vspace{4pt}\noindent
  \textbf{\color{orange!70!black}Osservazione.}\ 
}{%
  \par\vspace{4pt}
}

% ---------- Ambiente "Attenzione" (errore comune) ----------
\newenvironment{attenzione}{%
  \par\vspace{4pt}\noindent
  \begin{center}
  \begin{tabular}{|p{0.88\linewidth}|}
  \hline
  \rule{0pt}{2.3ex}\textbf{\color{red!70!black}Attenzione -- errore comune.}\ \ignorespaces
}{%
  \\[2pt]\hline
  \end{tabular}
  \end{center}
  \par\vspace{4pt}
}

% ---------- Box formula ----------
\newcommand{\formulabox}[1]{%
  \begin{center}
  \fbox{\parbox{0.9\linewidth}{\centering\vspace{2pt}#1\vspace{2pt}}}
  \end{center}
}

\setlist[itemize,1]{label=--}

\title{\textbf{Goniometria} \\ \large La misura degli angoli}
\author{Liceo Scientifico ``Galileo Galilei'' -- San Don\`a di Piave}
\date{Classe IV}

\begin{document}
\maketitle

\noindent
La \textbf{goniometria} \`e il ramo della matematica che studia gli angoli e le \textbf{funzioni goniometriche} associate ad essi. La \textbf{trigonometria} \`e invece il ramo della matematica che studia i procedimenti di calcolo che permettono di determinare gli elementi di un triangolo (lati e angoli) a partire dalla conoscenza di alcuni di essi.

\section{Misura degli angoli}

Per poter operare con gli angoli in modo algebrico \`e necessario disporre di un \emph{sistema di misura}, cio\`e associare a ogni angolo un numero reale che ne esprima l'ampiezza. I sistemi pi\`u utilizzati sono il \textbf{sistema sessagesimale} e il \textbf{sistema radiante}.

\subsection{Misura in gradi: il sistema sessagesimale}

Il \textbf{grado} \`e l'unit\`a di misura degli angoli nel sistema sessagesimale:
\formulabox{$1^\circ = \dfrac{1}{360}\ \text{dell'angolo giro}$}

Il sistema sessagesimale prende il nome dagli antichi \textbf{Babilonesi}, che utilizzavano il numero $60$ come base del loro sistema di numerazione. Le principali unit\`a di misura sono:
\[
1^\circ = 60' \qquad\qquad 1' = 60''
\]
dove $'$ indica i \textbf{primi} e $''$ indica i \textbf{secondi}.

\begin{esempio}
L'angolo
\[
32^\circ\ 13'\ 27''
\]
si legge: \emph{32 gradi, 13 primi e 27 secondi}.
\end{esempio}

\paragraph{Pregi e difetti del sistema sessagesimale.}
\begin{itemize}[nosep]
\item[\textcolor{green!50!black}{\textbf{+}}] \textbf{Pregi:} elevata precisione nelle misure angolari; ampio utilizzo in geometria, astronomia, navigazione e topografia.
\item[\textcolor{red!70!black}{\textbf{--}}] \textbf{Difetti:} le operazioni risultano pi\`u complesse rispetto al sistema decimale; \`e necessario effettuare conversioni tra gradi, primi e secondi.
\end{itemize}

\begin{attenzione}
Un valore decimale in gradi (come indicato, ad esempio, dalle calcolatrici) \textbf{non} coincide con la stessa cifra letta come gradi-primi-secondi. Il numero $37{,}25^\circ$ significa
\[
37^\circ + \frac{2^\circ}{10} + \frac{5^\circ}{100},
\]
cio\`e $37$ gradi e $0{,}25$ di grado (una frazione decimale di grado), e \textbf{non} $37$ gradi, $2$ primi e $5$ secondi:
\[
37{,}25^\circ \ \neq\ 37^\circ\,2'\,5''.
\]
Per convertire correttamente la parte decimale in primi e secondi si veda l'esempio seguente.
\end{attenzione}

\begin{esempio}[(conversione da gradi decimali a gradi-primi-secondi)]
Esprimere $32{,}275^\circ$ in gradi, primi e secondi.

\vspace{2pt}
\noindent La parte intera d\`a i gradi: $32^\circ$. Convertiamo la parte decimale in primi, moltiplicando per $60$:
\[
0{,}275^\circ \cdot 60 = 16{,}5' \ \Longrightarrow\ 16'\ \text{e}\ 0{,}5'\ \text{residui.}
\]
Convertiamo i primi residui in secondi, moltiplicando ancora per $60$:
\[
0{,}5' \cdot 60 = 30''.
\]
\[
\boxed{32{,}275^\circ = 32^\circ\,16'\,30''}
\]
\end{esempio}

\subsection{Misura in radianti}

Per semplificare i calcoli ed essere pi\`u pratici nell'ambito scientifico, in particolare in analisi matematica, si preferisce usare il \textbf{radiante}. Per definirlo consideriamo due circonferenze concentriche di raggio $r$ e $r'$ (con $r \neq r'$), aventi centro nel vertice di un angolo $\alpha$: l'angolo intercetta su di esse, rispettivamente, un arco $\ell$ e un arco $\ell'$.

\begin{center}
\begin{tikzpicture}[scale=1.15]
\draw (0,0) circle (1.1);
\draw (0,0) circle (2.0);
\draw[thick] (0,0) -- (2.3,0);
\draw[thick,blue] (0,0) -- (2.3*0.809,2.3*0.588);
\draw[very thick,red] (1.1,0) arc (0:36:1.1);
\draw[very thick,red] (2.0,0) arc (0:36:2.0);
\node at (0.55,0.10) {\small $\alpha$};
\draw (0.35,0) arc (0:36:0.35);
\node[red] at (1.55,0.42) {\small $\ell$};
\node[red] at (2.35,0.75) {\small $\ell'$};
\node at (1.35,-0.18) {\small $r$};
\node at (2.15,-0.18) {\small $r'$};
\end{tikzpicture}
\end{center}

Poich\'e gli angoli al centro sono direttamente proporzionali ai rispettivi archi, sulla prima circonferenza (circonferenza intera $=2\pi r$, angolo giro $=360^\circ$) vale la proporzione:
\[
\ell : \alpha^\circ = 2\pi r : 360^\circ
\qquad\Longrightarrow\qquad
\frac{\ell}{\alpha^\circ} = \frac{2\pi r}{360^\circ}
\qquad\Longrightarrow\qquad
\ell = \frac{2\pi r}{360^\circ}\cdot \alpha^\circ.
\]
Analogamente, sulla seconda circonferenza:
\[
\ell' : \alpha^\circ = 2\pi r' : 360^\circ
\qquad\Longrightarrow\qquad
\ell' = \frac{2\pi r'}{360^\circ}\cdot \alpha^\circ.
\]

Dividendo membro a membro le due proporzioni si ottiene $\ell : \ell' = r : r'$, cio\`e $\ell : r = \ell' : r'$, e quindi:
\formulabox{$\dfrac{\ell}{r} = \dfrac{\ell'}{r'}$}

Gli archi intercettati sono dunque \textbf{proporzionali} ai rispettivi raggi: il rapporto $\dfrac{\ell}{r}$ \textbf{non dipende dal raggio} della circonferenza scelta, ma \textbf{solo dall'ampiezza} dell'angolo $\alpha$. Infatti, da quanto ricavato sopra:
\[
\frac{\ell}{r} = 2\pi\,\frac{\alpha^\circ}{360^\circ} = \frac{\ell'}{r'}.
\]

\begin{center}
\fbox{\parbox{0.85\linewidth}{\centering\vspace{2pt}
\textbf{Definizione.} Si chiama \textbf{misura in radianti} di un angolo $\alpha$ il rapporto (costante, indipendente dal raggio) tra la lunghezza dell'arco intercettato su una qualunque circonferenza con centro nel vertice dell'angolo e la misura del raggio:
\[
\alpha_{\mathrm{rad}} = \frac{\ell}{r}.
\]
\vspace{2pt}}}
\end{center}

\begin{osservazione}
In particolare, se si sceglie $r=1$ (\textbf{circonferenza goniometrica}), risulta $\alpha_{\mathrm{rad}}=\ell$: la misura in radianti coincide numericamente con la lunghezza dell'arco intercettato. Questo giustifica la definizione equivalente: \emph{un radiante \`e l'ampiezza dell'angolo al centro che sottende un arco di lunghezza pari al raggio}. Inoltre, poich\'e \`e definito come rapporto fra due lunghezze, il radiante \`e un numero \textbf{puro} (adimensionale): non si scrive alcuna unit\`a di misura dopo il valore numerico.
\end{osservazione}

\subsection{Conversione tra gradi e radianti}

Ponendo $\alpha^\circ = 360^\circ$ (angolo giro) nella relazione $\dfrac{\ell}{r} = 2\pi\dfrac{\alpha^\circ}{360^\circ}$ si ottiene $\alpha_{\mathrm{rad}} = 2\pi$, coerentemente con il fatto che l'intera circonferenza ha lunghezza $2\pi r$. Si ricava cos\`i la proporzione fondamentale tra le due unit\`a di misura:
\formulabox{$360^\circ : 2\pi = \alpha^\circ : \alpha_{\mathrm{rad}} \qquad\Longleftrightarrow\qquad \alpha_{\mathrm{rad}} = \dfrac{\pi}{180^\circ}\cdot \alpha^\circ \qquad \alpha^\circ = \dfrac{180^\circ}{\pi}\cdot \alpha_{\mathrm{rad}}$}

\begin{esempio}
Convertire $150^\circ$ in radianti.
\[
\alpha_{\mathrm{rad}} = \frac{\pi}{180^\circ}\cdot 150^\circ = \frac{150\,\pi}{180} = \boxed{\frac{5}{6}\pi}
\]
\end{esempio}

\begin{esempio}
Convertire $\dfrac{3}{4}\pi$ radianti in gradi.
\[
\alpha^\circ = \frac{180^\circ}{\pi}\cdot \frac{3}{4}\pi = \boxed{135^\circ}
\]
\end{esempio}

\subsection{Tabella degli angoli notevoli}

\`E fondamentale memorizzare le misure, in gradi e in radianti, dei seguenti angoli notevoli, che ricorreranno continuamente nello studio delle funzioni goniometriche.

\begin{center}
\renewcommand{\arraystretch}{1.6}
\begin{tabular}{|c|c|c|c|c|c|c|c|c|c|}
\hline
\textbf{Gradi} & $0^\circ$ & $30^\circ$ & $45^\circ$ & $60^\circ$ & $90^\circ$ & $120^\circ$ & $135^\circ$ & $150^\circ$ & $180^\circ$ \\
\hline
\textbf{Radianti} & $0$ & $\dfrac{\pi}{6}$ & $\dfrac{\pi}{4}$ & $\dfrac{\pi}{3}$ & $\dfrac{\pi}{2}$ & $\dfrac{2\pi}{3}$ & $\dfrac{3\pi}{4}$ & $\dfrac{5\pi}{6}$ & $\pi$ \\
\hline
\end{tabular}
\end{center}

\begin{center}
\renewcommand{\arraystretch}{1.6}
\begin{tabular}{|c|c|c|c|c|c|c|c|c|}
\hline
\textbf{Gradi} & $210^\circ$ & $225^\circ$ & $240^\circ$ & $270^\circ$ & $300^\circ$ & $315^\circ$ & $330^\circ$ & $360^\circ$ \\
\hline
\textbf{Radianti} & $\dfrac{7\pi}{6}$ & $\dfrac{5\pi}{4}$ & $\dfrac{4\pi}{3}$ & $\dfrac{3\pi}{2}$ & $\dfrac{5\pi}{3}$ & $\dfrac{7\pi}{4}$ & $\dfrac{11\pi}{6}$ & $2\pi$ \\
\hline
\end{tabular}
\end{center}

\begin{center}
\begin{tikzpicture}[scale=1.7]
\draw[->] (-1.4,0) -- (1.4,0);
\draw[->] (0,-1.4) -- (0,1.4);
\draw (0,0) circle (1);
\foreach \ang/\lab in {0/0,30/{\pi/6},45/{\pi/4},60/{\pi/3},90/{\pi/2},120/{2\pi/3},135/{3\pi/4},150/{5\pi/6},180/\pi,210/{7\pi/6},225/{5\pi/4},240/{4\pi/3},270/{3\pi/2},300/{5\pi/3},315/{7\pi/4},330/{11\pi/6}}{
  \draw[gray] (0,0) -- (\ang:1);
  \filldraw[blue] (\ang:1) circle (0.4pt);
  \node[font=\tiny] at (\ang:1.22) {$\lab$};
}
\end{tikzpicture}
\end{center}

\section{Lunghezza di un arco di circonferenza}

Dalla definizione di radiante, $\alpha_{\mathrm{rad}} = \dfrac{\ell}{r}$, si ricava immediatamente la formula per la lunghezza dell'arco $\ell$ sotteso da un angolo al centro $\alpha$ (espresso in radianti) su una circonferenza di raggio $r$:
\formulabox{$\ell = \alpha \cdot r$}

\begin{esempio}
Calcolare la lunghezza dell'arco sotteso da un angolo al centro di $60^\circ$ su una circonferenza di raggio $r = 9\ \text{cm}$.

\vspace{2pt}
\noindent Convertiamo l'angolo in radianti: $60^\circ = \dfrac{\pi}{3}$. Quindi:
\[
\ell = \frac{\pi}{3}\cdot 9 = \boxed{3\pi\ \text{cm} \approx 9{,}42\ \text{cm}}
\]
\end{esempio}

\section{Area del settore circolare}

Un settore circolare \`e la parte di cerchio delimitata da un angolo al centro $\alpha$ e dall'arco corrispondente. Poich\'e l'area del settore \`e direttamente proporzionale all'ampiezza dell'angolo al centro, si pu\`o impostare la proporzione con l'angolo giro (a cui corrisponde l'intero cerchio, di area $\pi r^2$):
\[
A : \pi r^2 = \alpha : 2\pi
\]
da cui, semplificando:
\formulabox{$A = \dfrac{1}{2}\,\alpha\, r^2 \qquad (\alpha\ \text{in radianti})$}

\begin{osservazione}
Poich\'e $\ell = \alpha r$, si pu\`o riscrivere la formula anche come $A = \dfrac{1}{2}\,\ell\, r$, in perfetta analogia con la formula dell'area del triangolo.
\end{osservazione}

\begin{esempio}
Calcolare l'area del settore circolare di raggio $r = 6\ \text{cm}$ e angolo al centro $\alpha = \dfrac{2}{3}\pi$.
\[
A = \frac{1}{2}\cdot \frac{2\pi}{3}\cdot 6^2 = \frac{1}{2}\cdot \frac{2\pi}{3}\cdot 36 = \boxed{12\pi\ \text{cm}^2 \approx 37{,}70\ \text{cm}^2}
\]
\end{esempio}

\section{Angoli orientati}

Fino a questo punto abbiamo considerato l'ampiezza di un angolo come un valore non negativo. In goniometria, per\`o, \`e fondamentale associare a ogni angolo anche un \textbf{verso di rotazione}, distinguendo il lato origine dal lato termine.

\begin{center}
\fbox{\parbox{0.85\linewidth}{\centering\vspace{2pt}
Un angolo si dice \textbf{orientato} quando si stabilisce quale dei due lati \`e il \emph{lato origine} e quale il \emph{lato termine}, individuando cos\`i un verso di rotazione dal primo al secondo.\vspace{2pt}}}
\end{center}

Per convenzione:
\begin{itemize}[nosep]
\item un angolo si dice \textbf{positivo} se la rotazione dal lato origine al lato termine avviene in senso \textbf{antiorario};
\item un angolo si dice \textbf{negativo} se la rotazione avviene in senso \textbf{orario}.
\end{itemize}

\begin{center}
\begin{tikzpicture}[scale=1.1]
\begin{scope}
\draw[->] (0,0) -- (1.6,0) node[right] {\small lato origine};
\draw[->] (0,0) -- (0.98,1.26) node[above right] {\small lato termine};
\draw[->,thick,blue] (0.7,0) arc (0:52:0.7);
\node[blue] at (0.55,0.45) {\small $+\alpha$};
\node at (0.8,-0.7) {\small angolo \textbf{positivo} (antiorario)};
\end{scope}
\begin{scope}[xshift=6.5cm]
\draw[->] (0,0) -- (1.6,0) node[right] {\small lato origine};
\draw[->] (0,0) -- (0.98,-1.26) node[below right] {\small lato termine};
\draw[->,thick,red] (0.7,0) arc (0:-52:0.7);
\node[red] at (0.55,-0.45) {\small $-\alpha$};
\node at (0.8,-1.9) {\small angolo \textbf{negativo} (orario)};
\end{scope}
\end{tikzpicture}
\end{center}

\begin{osservazione}
Grazie all'orientamento, due angoli che differiscono per un multiplo intero dell'angolo giro (cio\`e $\alpha$ e $\alpha + k\cdot 360^\circ$, con $k \in \mathbb{Z}$, oppure $\alpha + k\cdot 2\pi$ in radianti) individuano \emph{la stessa posizione finale} sulla circonferenza, pur essendo considerati angoli distinti dal punto di vista della rotazione compiuta. Questa idea sar\`a alla base della definizione delle funzioni goniometriche periodiche.
\end{osservazione}

\section{Esercizi di riepilogo (misura degli angoli)}

\begin{enumerate}[label=\arabic*)]
\item Converti in gradi, primi e secondi: $18{,}625^\circ$.
\item Spiega perch\'e $54{,}5^\circ \neq 54^\circ\,5'\,0''$ e trova la scrittura corretta in gradi-primi-secondi.
\item Converti in radianti: $75^\circ$, $\ 252^\circ$.
\item Converti in gradi: $\dfrac{7}{12}\pi$, $\ \dfrac{5}{9}\pi$.
\item Calcola la lunghezza dell'arco sotteso da un angolo al centro di $45^\circ$ su una circonferenza di raggio $r = 12\ \text{cm}$.
\item Calcola l'area del settore circolare di raggio $r = 5\ \text{cm}$ e angolo al centro $\alpha = \dfrac{3}{4}\pi$.
\item Un settore circolare ha area $A = 20\pi\ \text{cm}^2$ ed \`e generato da un angolo al centro di $72^\circ$. Determina il raggio della circonferenza.
\end{enumerate}

\section{La circonferenza goniometrica}

Per definire in modo unitario le funzioni goniometriche di \emph{qualunque} angolo (non solo di angoli acuti, come nella trigonometria del triangolo rettangolo) si introduce la \textbf{circonferenza goniometrica}.

\begin{center}
\fbox{\parbox{0.85\linewidth}{\centering\vspace{2pt}
\textbf{Definizione.} Si chiama \textbf{circonferenza goniometrica} la circonferenza con centro nell'origine $O$ di un sistema di riferimento cartesiano ortogonale e raggio $r=1$.\vspace{2pt}}}
\end{center}

Un angolo orientato $\alpha$ si rappresenta sulla circonferenza goniometrica facendo coincidere il \emph{lato origine} con il semiasse positivo delle ascisse e il vertice con $O$: il \emph{lato termine} interseca la circonferenza in un unico punto $P$, detto \textbf{punto associato} all'angolo $\alpha$.

\begin{center}
\begin{tikzpicture}[scale=1.7]
\draw[->] (-1.5,0) -- (1.5,0) node[right] {\small $x$};
\draw[->] (0,-1.5) -- (0,1.5) node[above] {\small $y$};
\draw (0,0) circle (1);
\draw[thick,blue,->] (0,0) -- (40:1) node[pos=1.15] {\small $P$};
\draw[dashed] (40:1) -- (40:1 |- 0,0) node[below] {\small $H$};
\draw[dashed] (40:1) -- (0,{sin(40)});
\draw (0.28,0) arc (0:40:0.28);
\node at (0.45,0.13) {\small $\alpha$};
\node[below left] at (0,0) {\small $O$};
\node[below] at (1,0) {\small $1$};
\end{tikzpicture}
\end{center}

\section{Seno e coseno di un angolo}

Dato un angolo orientato $\alpha$ e il corrispondente punto $P(x_P,y_P)$ sulla circonferenza goniometrica, si definiscono:

\formulabox{$\cos\alpha = x_P \qquad\qquad \sin\alpha = y_P$}

cio\`e il \textbf{coseno} di $\alpha$ \`e l'ascissa di $P$ e il \textbf{seno} di $\alpha$ \`e l'ordinata di $P$.

\begin{osservazione}
Per un angolo acuto $\alpha$ di un triangolo rettangolo, questa definizione \`e coerente con quella nota dalla geometria:
\[
\sin\alpha = \frac{\text{cateto opposto}}{\text{ipotenusa}}, \qquad \cos\alpha = \frac{\text{cateto adiacente}}{\text{ipotenusa}},
\]
poich\'e sulla circonferenza goniometrica l'ipotenusa del triangolo $OHP$ \`e proprio il raggio $r=1$.
\end{osservazione}

\paragraph{Propriet\`a fondamentali.} Poich\'e $P$ appartiene alla circonferenza di raggio $1$:
\begin{itemize}[nosep]
\item $-1 \le \sin\alpha \le 1$ \ e \ $-1 \le \cos\alpha \le 1$ \ per ogni $\alpha$;
\item $\mathrm{Dom}(\sin) = \mathrm{Dom}(\cos) = \mathbb{R}$: seno e coseno sono definiti per ogni valore reale dell'angolo;
\item \textbf{periodicit\`a}: $\sin(\alpha + 2k\pi) = \sin\alpha$ \ e \ $\cos(\alpha + 2k\pi) = \cos\alpha$, per ogni $k\in\mathbb{Z}$ (il punto $P$ torna nella stessa posizione dopo un giro completo);
\item il coseno \`e una funzione \textbf{pari}: $\cos(-\alpha) = \cos\alpha$;
\item il seno \`e una funzione \textbf{dispari}: $\sin(-\alpha) = -\sin\alpha$.
\end{itemize}

\section{Tangente di un angolo}

Consideriamo la retta $t$ tangente alla circonferenza goniometrica nel punto $(1,0)$, parallela all'asse $y$ (detta \textbf{asse delle tangenti}). Prolungando il raggio $OP$ fino a incontrare $t$ in un punto $T$, l'ordinata di $T$ definisce la \textbf{tangente} dell'angolo $\alpha$:

\formulabox{$\tan\alpha = \dfrac{\sin\alpha}{\cos\alpha} \qquad (\cos\alpha \neq 0)$}

\begin{center}
\begin{tikzpicture}[scale=1.7]
\draw[->] (-1.5,0) -- (2.0,0) node[right] {\small $x$};
\draw[->] (0,-1.5) -- (0,1.6) node[above] {\small $y$};
\draw (0,0) circle (1);
\draw[dashed] (1,-1.4) -- (1,1.4) node[above] {\small $t$};
\draw[thick,blue,->] (0,0) -- (50:1.55) node[pos=0.66,above right] {\small $P$};
\filldraw[black] (50:1) circle (0.4pt);
\node[blue] at (1,1.19) {\small $T$};
\filldraw[blue] (1,1.19) circle (0.4pt);
\draw (0.28,0) arc (0:50:0.28);
\node at (0.48,0.15) {\small $\alpha$};
\node[below left] at (0,0) {\small $O$};
\end{tikzpicture}
\end{center}

\begin{osservazione}
La tangente \textbf{non} \`e definita quando $\cos\alpha = 0$, cio\`e per
\[
\alpha = \frac{\pi}{2} + k\pi, \qquad k \in \mathbb{Z}
\]
(in questi casi la retta $OP$ \`e parallela all'asse delle tangenti e non la incontra mai). Inoltre $\tan\alpha$ \`e periodica di periodo $\pi$ (non $2\pi$ come seno e coseno):
\[
\tan(\alpha + k\pi) = \tan\alpha, \qquad k \in \mathbb{Z}.
\]
\end{osservazione}

\section{Relazione fondamentale della goniometria}

Poich\'e $P(\cos\alpha,\sin\alpha)$ appartiene alla circonferenza goniometrica $x^2+y^2=1$, vale per ogni angolo $\alpha$ la relazione:
\formulabox{$\sin^2\alpha + \cos^2\alpha = 1$}

Da questa relazione, nota $\sin\alpha$ (oppure $\cos\alpha$), si ricava il \textbf{valore assoluto} dell'altra funzione, poich\'e la radice quadrata restituisce sempre un valore non negativo:
\[
|\cos\alpha| = \sqrt{1-\sin^2\alpha}, \qquad\qquad |\sin\alpha| = \sqrt{1-\cos^2\alpha}.
\]
Il segno effettivo di $\cos\alpha$ (o di $\sin\alpha$) va poi stabilito separatamente, osservando il quadrante in cui si trova $\alpha$ (si veda la tabella dei segni, Sezione~\ref{sec:segno-quadranti}):
\[
\cos\alpha = \pm|\cos\alpha|, \qquad\qquad \sin\alpha = \pm|\sin\alpha|.
\]

Dividendo la relazione fondamentale per $\cos^2\alpha$ (con $\cos\alpha\neq 0$) si ottiene inoltre una relazione utile tra tangente e coseno:
\formulabox{$1+\tan^2\alpha = \dfrac{1}{\cos^2\alpha}$}

\begin{esempio}
Sapendo che $\sin\alpha = \dfrac{3}{5}$ e che $\alpha$ \`e un angolo ottuso (secondo quadrante), calcolare $\cos\alpha$ e $\tan\alpha$.

\vspace{2pt}
\noindent Dalla relazione fondamentale:
\[
\cos^2\alpha = 1-\sin^2\alpha = 1-\frac{9}{25}=\frac{16}{25} \ \Longrightarrow\ \cos\alpha = \pm\frac{4}{5}.
\]
Poich\'e nel secondo quadrante il coseno \`e negativo:
\[
\boxed{\cos\alpha = -\frac{4}{5}} \qquad\qquad \tan\alpha = \frac{\sin\alpha}{\cos\alpha} = \frac{3/5}{-4/5} = \boxed{-\frac{3}{4}}
\]
\end{esempio}

\section{Valori delle funzioni goniometriche degli angoli notevoli}

\begin{center}
\renewcommand{\arraystretch}{2.0}
\begin{tabular}{|c|c|c|c|c|c|}
\hline
$\alpha$ & $0$ & $\dfrac{\pi}{6}$ & $\dfrac{\pi}{4}$ & $\dfrac{\pi}{3}$ & $\dfrac{\pi}{2}$ \\
\hline
$\sin\alpha$ & $0$ & $\dfrac{1}{2}$ & $\dfrac{\sqrt{2}}{2}$ & $\dfrac{\sqrt{3}}{2}$ & $1$ \\
\hline
$\cos\alpha$ & $1$ & $\dfrac{\sqrt{3}}{2}$ & $\dfrac{\sqrt{2}}{2}$ & $\dfrac{1}{2}$ & $0$ \\
\hline
$\tan\alpha$ & $0$ & $\dfrac{\sqrt{3}}{3}$ & $1$ & $\sqrt{3}$ & \text{n.e.} \\
\hline
\end{tabular}
\end{center}

\begin{osservazione}
``n.e.'' significa \emph{non esiste}: la tangente non \`e definita per $\alpha=\dfrac{\pi}{2}$, poich\'e $\cos\dfrac{\pi}{2}=0$. I valori delle funzioni goniometriche per gli angoli dei restanti quadranti si ricavano da questi tramite le formule degli \textbf{angoli associati} (Sezione~\ref{sec:angoli-associati}).
\end{osservazione}

\section{Segno delle funzioni goniometriche nei quadranti}
\label{sec:segno-quadranti}

Il segno di $\sin\alpha$, $\cos\alpha$ e $\tan\alpha$ dipende dal quadrante in cui si trova il punto $P$ associato all'angolo $\alpha$:

\begin{center}
\renewcommand{\arraystretch}{1.6}
\begin{tabular}{|c|c|c|c|c|}
\hline
 & \textbf{I quadrante} & \textbf{II quadrante} & \textbf{III quadrante} & \textbf{IV quadrante} \\
 & $\left(0,\dfrac{\pi}{2}\right)$ & $\left(\dfrac{\pi}{2},\pi\right)$ & $\left(\pi,\dfrac{3\pi}{2}\right)$ & $\left(\dfrac{3\pi}{2},2\pi\right)$ \\
\hline
$\sin\alpha$ & $+$ & $+$ & $-$ & $-$ \\
\hline
$\cos\alpha$ & $+$ & $-$ & $-$ & $+$ \\
\hline
$\tan\alpha$ & $+$ & $-$ & $+$ & $-$ \\
\hline
\end{tabular}
\end{center}

\section{Angoli associati}
\label{sec:angoli-associati}

Si dicono \textbf{angoli associati} due o pi\`u angoli i cui punti sulla circonferenza goniometrica sono legati da una relazione di simmetria; le loro funzioni goniometriche sono quindi legate tra loro. Le relazioni pi\`u utilizzate sono le seguenti (con $\alpha$ angolo qualunque):

\begin{center}
\renewcommand{\arraystretch}{1.8}
\begin{tabular}{|l|c|c|c|}
\hline
\textbf{Relazione} & $\sin$ & $\cos$ & $\tan$ \\
\hline
Angoli \textbf{opposti} $-\alpha$ & $\sin(-\alpha)=-\sin\alpha$ & $\cos(-\alpha)=\cos\alpha$ & $\tan(-\alpha)=-\tan\alpha$ \\
\hline
Angoli \textbf{supplementari} $\pi-\alpha$ & $\sin(\pi-\alpha)=\sin\alpha$ & $\cos(\pi-\alpha)=-\cos\alpha$ & $\tan(\pi-\alpha)=-\tan\alpha$ \\
\hline
Angoli che differiscono di $\pi$ & $\sin(\pi+\alpha)=-\sin\alpha$ & $\cos(\pi+\alpha)=-\cos\alpha$ & $\tan(\pi+\alpha)=\tan\alpha$ \\
\hline
Angoli \textbf{esplementari} $2\pi-\alpha$ & $\sin(2\pi-\alpha)=-\sin\alpha$ & $\cos(2\pi-\alpha)=\cos\alpha$ & $\tan(2\pi-\alpha)=-\tan\alpha$ \\
\hline
Angoli \textbf{complementari} $\dfrac{\pi}{2}-\alpha$ & $\sin\!\left(\dfrac{\pi}{2}-\alpha\right)=\cos\alpha$ & $\cos\!\left(\dfrac{\pi}{2}-\alpha\right)=\sin\alpha$ & -- \\
\hline
\end{tabular}
\end{center}

\begin{esempio}
Calcolare $\sin 150^\circ$ usando gli angoli associati.

\vspace{2pt}
\noindent Poich\'e $150^\circ = 180^\circ - 30^\circ$, si tratta di angoli supplementari:
\[
\sin 150^\circ = \sin(180^\circ - 30^\circ) = \sin 30^\circ = \boxed{\frac{1}{2}}
\]
\end{esempio}

\section{Grafici delle funzioni goniometriche}

\subsection{La sinusoide: $y=\sin x$}
\begin{center}
\begin{tikzpicture}[scale=1.0]
\draw[->] (-0.5,0) -- (7,0) node[right] {\small $x$};
\draw[->] (0,-1.4) -- (0,1.4) node[above] {\small $y$};
\draw[thick,blue,domain=0:6.4,samples=100] plot (\x,{sin(\x r)});
\foreach \x/\lab in {1.5708/{\pi/2},3.1416/\pi,4.7124/{3\pi/2},6.2832/{2\pi}}{
  \draw[dashed,gray] (\x,0) -- (\x,{sin(\x r)});
  \node[below] at (\x,-0.05) {\small $\lab$};
}
\draw[dashed] (-0.5,1) -- (7,1); \node[left] at (0,1) {\small $1$};
\draw[dashed] (-0.5,-1) -- (7,-1); \node[left] at (0,-1) {\small $-1$};
\end{tikzpicture}
\end{center}
Dominio $\mathbb{R}$, codominio $[-1,1]$, periodo $2\pi$, funzione dispari.

\subsection{La cosinusoide: $y=\cos x$}
\begin{center}
\begin{tikzpicture}[scale=1.0]
\draw[->] (-0.5,0) -- (7,0) node[right] {\small $x$};
\draw[->] (0,-1.4) -- (0,1.4) node[above] {\small $y$};
\draw[thick,red,domain=0:6.4,samples=100] plot (\x,{cos(\x r)});
\foreach \x/\lab in {1.5708/{\pi/2},3.1416/\pi,4.7124/{3\pi/2},6.2832/{2\pi}}{
  \draw[dashed,gray] (\x,0) -- (\x,{cos(\x r)});
  \node[below] at (\x,-0.05) {\small $\lab$};
}
\draw[dashed] (-0.5,1) -- (7,1); \node[left] at (0,1) {\small $1$};
\draw[dashed] (-0.5,-1) -- (7,-1); \node[left] at (0,-1) {\small $-1$};
\end{tikzpicture}
\end{center}
Dominio $\mathbb{R}$, codominio $[-1,1]$, periodo $2\pi$, funzione pari. \`E la sinusoide traslata di $\dfrac{\pi}{2}$ verso sinistra: $\cos x = \sin\!\left(x+\dfrac{\pi}{2}\right)$.

\subsection{La tangentoide: $y=\tan x$}
\begin{center}
\begin{tikzpicture}[scale=1.0]
\draw[->] (-0.3,0) -- (6.6,0) node[right] {\small $x$};
\draw[->] (0,-2.2) -- (0,2.2) node[above] {\small $y$};
\begin{scope}
\clip (-0.3,-2.1) rectangle (6.6,2.1);
\draw[thick,green!50!black,domain=-1.45:1.45,samples=150] plot (\x+1.5708,{tan(\x r)});
\draw[thick,green!50!black,domain=-1.45:1.45,samples=150] plot (\x+4.7124,{tan(\x r)});
\end{scope}
\draw[dashed,gray] (1.5708,-2.2) -- (1.5708,2.2); \node[below] at (1.5708,-2.25) {\small $\pi/2$};
\draw[dashed,gray] (4.7124,-2.2) -- (4.7124,2.2); \node[below] at (4.7124,-2.25) {\small $3\pi/2$};
\node[below] at (3.1416,-0.05) {\small $\pi$};
\draw[dashed,gray] (3.1416,0) -- (3.1416,0.05);
\end{tikzpicture}
\end{center}
Dominio $\mathbb{R}\setminus\left\{\dfrac{\pi}{2}+k\pi\right\}$, codominio $\mathbb{R}$, periodo $\pi$, funzione dispari; asintoti verticali nei punti esclusi dal dominio.

\section{Esercizi di riepilogo (seno, coseno, tangente)}

\begin{enumerate}[label=\arabic*)]
\item Utilizzando la circonferenza goniometrica, determina $\sin\alpha$, $\cos\alpha$, $\tan\alpha$ per $\alpha = \dfrac{5\pi}{6}$.
\item Sapendo che $\cos\alpha = -\dfrac{5}{13}$ e che $\alpha$ \`e un angolo del terzo quadrante, calcola $\sin\alpha$ e $\tan\alpha$.
\item Verifica, usando le formule degli angoli associati, che $\cos 210^\circ = -\cos 30^\circ$.
\item Calcola $\sin 300^\circ$ e $\cos 300^\circ$ riconducendoti a un angolo del primo quadrante.
\item Stabilisci il segno di $\sin 250^\circ$, $\cos 250^\circ$ e $\tan 250^\circ$ senza calcolarne il valore, indicando il quadrante di appartenenza.
\item Sapendo che $\tan\alpha = \dfrac{4}{3}$ e che $\alpha$ \`e un angolo del terzo quadrante, calcola $\sin\alpha$ e $\cos\alpha$.
\end{enumerate}

\end{document}